\chapter{Formelzeichen und Abkürzungen}\label{chapter:symbols}


\subsubsection{Formelzeichen -- Lateinische Buchstaben}

\begin{tabularx}{\textwidth}{cXc}
Variable & Bezeichnung & Dimension \\ \\
  $a$              & Temperaturleitfähigkeit                & m$^2$/s \\
  $A$              & Flammenoberfläche                      & m$^2$ \\
  $A_i$            & präexponentieller Faktor der Reaktion $i$ & -- \\
  $c_p$            & spezifische isobare Wärmekapazität     & J/(kg\,K) \\
  $D$              & Diffusionskoeffizient                  & m$^2$/s \\
  $E_\mathrm{a}$   & Aktivierungsenergie                    & J/mol \\
  $g_\mathrm{c}$   & kritischer Geschwindigkeitsgradient    & 1/s \\
  $h$              & spezifische Enthalpie                  & J/kg \\
  $k_\mathrm{f}$   & Geschwindigkeitskoeffizient (Hinreaktion) & -- \\
  $l_\mathrm{t}$   & integrales turbulentes Längenmaß       & m \\
  $\dot{m}$        & Massenstromdichte                      & kg/(m$^2$\,s) \\
  $p$              & Druck                                  & Pa \\
  $q_i$            & Netto-Reaktionsrate der Reaktion $i$   & mol/(m$^3$\,s) \\
  $R_\mathrm{m}$   & universelle Gaskonstante               & J/(mol\,K) \\
  $s_\mathrm{L}$   & laminare Flammengeschwindigkeit        & m/s \\
  $s_\mathrm{T}$   & turbulente Flammengeschwindigkeit      & m/s \\
  $T$              & Temperatur                             & K \\
  $T_\mathrm{ad}$  & adiabate Flammentemperatur             & K \\
  $u$              & Strömungsgeschwindigkeit               & m/s \\
  $u^{\prime}$     & turbulente Schwankungsgeschwindigkeit  & m/s \\
  $V_k$            & Diffusionsgeschwindigkeit der Spezies $k$ & m/s \\
  $\dot{V}$        & Volumenstrom                           & m$^3$/s \\
  $W_k$            & molare Masse der Spezies $k$           & kg/mol \\
  $x$              & Ortskoordinate                         & m \\
  $Y_k$            & Massenbruch der Spezies $k$            & -- \\
\end{tabularx}


\subsubsection{Formelzeichen -- Griechische Buchstaben}

\begin{tabularx}{\textwidth}{cXc}
Variable & Bezeichnung & Dimension \\ \\
  $\alpha$              & Temperaturexponent                     & -- \\
  $\beta$               & Druckexponent                          & -- \\
  $\delta_\mathrm{L}$   & laminare Flammendicke                  & m \\
  $\delta_\mathrm{q}$   & Löschabstand                           & m \\
  $\eta$                & Kolmogorov-Längenmaß                   & m \\
  $\kappa$              & Streckungsrate                         & 1/s \\
  $\lambda$             & Wärmeleitfähigkeit                     & W/(m\,K) \\
  $\lambda$             & Luftzahl                               & -- \\
  $\mathcal{L}$         & Markstein-Länge                        & m \\
  $\nu^{\prime}$        & stöchiometrischer Koeffizient (Edukte) & -- \\
  $\nu^{\prime\prime}$  & stöchiometrischer Koeffizient (Produkte) & -- \\
  $\rho$                & Dichte                                 & kg/m$^3$ \\
  $\tau_\mathrm{c}$     & chemisches Zeitmaß                     & s \\
  $\tau_\mathrm{t}$     & turbulentes Zeitmaß                    & s \\
  $\varphi$             & Äquivalenzverhältnis                   & -- \\
  $\dot{\omega}_k$      & molare Bildungsrate der Spezies $k$    & mol/(m$^3$\,s) \\
\end{tabularx}


\subsubsection{Dimensionslose Kennzahlen}

\begin{tabularx}{\textwidth}{cXc}
Kennzahl & Bezeichnung & Definition \\ \\
  $\mathrm{Da}$ & Damköhler-Zahl  & $\tau_\mathrm{t}/\tau_\mathrm{c}$ \\
  $\mathrm{Ka}$ & Karlovitz-Zahl  & $\tau_\mathrm{c}/\tau_\eta$ \\
  $\mathrm{Le}$ & Lewis-Zahl      & $a/D$ \\
  $\mathrm{Ma}$ & Markstein-Zahl  & $\mathcal{L}/\delta_\mathrm{L}$ \\
  $\mathrm{Pe}$ & Péclet-Zahl     & $g_\mathrm{c}\delta_\mathrm{L}/s_\mathrm{L}$ \\
  $\mathrm{Re}$ & Reynolds-Zahl   & $u\,d/\nu$ \\
\end{tabularx}


\subsubsection{Abkürzungen}

\begin{tabularx}{\textwidth}{lX}
BLF   & Boundary Layer Flashback (Grenzschichtrückschlag) \\
CFD   & Computational Fluid Dynamics \\
CIVB  & Combustion Induced Vortex Breakdown \\
LTT   & Lehrstuhl für Technische Thermodynamik \\
MFC   & Mass Flow Controller (Massendurchflussregler) \\
PIV   & Particle Image Velocimetry \\
PLIF  & Planar Laser Induced Fluorescence \\
\end{tabularx}


\subsubsection{Indizes}

\begin{tabularx}{\textwidth}{cX}
  $\mathrm{ad}$ & adiabat \\
  $\mathrm{b}$  & verbrannt (burnt) \\
  $\mathrm{c}$  & kritisch bzw.\ chemisch \\
  $k$           & Laufindex der Spezies \\
  $\mathrm{L}$  & laminar \\
  $\mathrm{st}$ & stöchiometrisch \\
  $\mathrm{T}$  & turbulent \\
  $\mathrm{u}$  & unverbrannt (unburnt) \\
  $0$           & Referenzzustand \\
\end{tabularx}
